\documentclass[../../book]{subfiles}

\begin{document}

\subsection{Interactive Oracle Proofs of Proximity (IOPP)}

Recall the idea of an \emph{interactive oracle proof} (IOP), which we already met in \Cref{section:plonk}. % TODO: point to where IOPs were introduced
As in any interactive proof, the prover and the verifier exchange messages over several rounds, and the verifier's messages are random challenges.
The difference is that the prover's messages are \emph{oracles}: long strings that the verifier does not read in full, but only queries at a few positions of its choice (see \Cref{fig:iop}).
This is what allows the verifier to run in time much smaller than the total length of the prover's messages.
In practice, the prover commits to each oracle with a Merkle tree (\Cref{section:commitment-schemes}), and every query is answered with an opening.

\begin{figure}[h!]
    \centering
    \scalebox{0.8}{%
    \begin{tikzpicture}[
        cell/.style={draw=oc-gray-7, minimum size=0.36cm, inner sep=0pt},
        send/.style={-{Stealth[length=3mm]}, line width=0.4mm},
    ]
        \node[inner sep=0pt, align=center] (prover) {\includegraphics[width=1.25cm]{contents/images/common/prover.png}\\Prover $\mathcal{P}(\mathbbm{x}, \mathbbm{w})$};
        \node[inner sep=0pt, align=center, right=5cm of prover] (verifier) {\includegraphics[width=1.25cm]{contents/images/common/verifier.png}\\Verifier $\mathcal{V}(\mathbbm{x})$};

        \draw [dashed,line width=0.3mm] ([yshift=-0.5cm]prover.south) -- ([yshift=-9.1cm]prover.south);
        \draw [dashed,line width=0.3mm] ([yshift=-0.5cm]verifier.south) -- ([yshift=-9.1cm]verifier.south);

        % an oracle sent by the prover: #1 = yshift, #2 = name, #3 = queried cells
        \newcommand{\sendoracle}[3]{
            \draw[send] ([yshift=#1]prover.south) coordinate (l) -- (l-|verifier) coordinate[midway] (m);
            \foreach \i in {0,...,9} {
                \node[cell, fill=oc-blue-1] at ([xshift=-1.1cm+0.36cm*\i, yshift=0.35cm]m) {};
            }
            \foreach \i in {#3} {
                \node[cell, fill=oc-orange-3, draw=oc-orange-8, thick] at ([xshift=-1.1cm+0.36cm*\i, yshift=0.35cm]m) {};
            }
            \node[left] at ([xshift=-1.3cm, yshift=0.35cm]m) {#2 $=$};
        }
        % a challenge sent by the verifier: #1 = arrow yshift, #2 = name
        \newcommand{\sendchallenge}[2]{
            \draw[send] ([yshift=#1]verifier.south) coordinate (l) -- (l-|prover) node[midway, above=0mm]{Sample challenge #2};
        }

        \sendoracle{-1.6cm}{$\pi_0$}{2,7}
        \sendchallenge{-2.8cm}{$\alpha_1$}
        \sendoracle{-4.2cm}{$\pi_1$}{4}

        \path (prover.south) -- (verifier.south) coordinate[midway] (c);
        \node at ([yshift=-4.8cm]c) {$\vdots$};

        \sendchallenge{-6.0cm}{$\alpha_r$}
        \sendoracle{-7.4cm}{$\pi_r$}{0,5,8}

        \node[align=center,fill=white!5,thick,below=7.8cm of verifier]{
        \noindent\rule{3.5cm}{0.8pt}\\
        Query a few cells of each $\pi_i$\\
        Output $\mathsf{accept}/\mathsf{reject}$ \\
        \noindent\rule{3.5cm}{0.8pt}};
    \end{tikzpicture}%
    }

    \caption{An interactive oracle proof between prover $\mathcal{P}$ and verifier $\mathcal{V}$ with $r$ rounds of interaction. The prover's messages $\pi_0, \dots, \pi_r$ are long strings (oracles), and the verifier's messages $\alpha_1, \dots, \alpha_r$ are short random challenges. The verifier does not read the oracles in full: it queries only a few cells of each (marked in orange) and decides based on the answers.}
    \label{fig:iop-definition}
\end{figure}

In an IOP, the verifier reads the statement $\mathbbm{x}$ in full.
For the proximity problem this is not enough: the object we want to check, a word $\mathbf{w}$, is itself too long to read.
We therefore let a part of the statement be an oracle as well.

\begin{definition}[IOP of Proximity]\label{def:iopp}
    Let $\mathcal{R} = \{(\mathbbm{x}, \mathbbm{y}, \mathbbm{w})\}$ be a relation, where $\mathbbm{y}$ is a string of large length (say, $n$).
    An \emph{IOP of proximity} (IOPP) for $\mathcal{R}$ is an IOP in which the prover receives $(\mathbbm{x}, \mathbbm{y}, \mathbbm{w})$, while the verifier receives $\mathbbm{x}$ and only \emph{oracle access} to $\mathbbm{y}$: it may query cells of $\mathbbm{y}$ in the same way as cells of the oracles $\pi_0, \dots, \pi_r$.
    At the end, the verifier outputs
    \begin{equation*}
        \mathcal{V}^{\mathbbm{y}, \pi_0, \dots, \pi_r}(\mathbbm{x}, \alpha_1, \dots, \alpha_r) \to \{0,1\}.
    \end{equation*}
\end{definition}

\begin{figure}[h!]
    \centering
    \scalebox{0.8}{%
    \begin{tikzpicture}[
        cell/.style={draw=oc-gray-7, minimum size=0.36cm, inner sep=0pt},
        send/.style={-{Stealth[length=3mm]}, line width=0.4mm},
    ]
        \node[inner sep=0pt, align=center] (prover) {\includegraphics[width=1.25cm]{contents/images/common/prover.png}\\Prover $\mathcal{P}(\mathbbm{x}, \textcolor{oc-red-8}{\mathbbm{y}}, \mathbbm{w})$};
        \node[inner sep=0pt, align=center, right=5cm of prover] (verifier) {\includegraphics[width=1.25cm]{contents/images/common/verifier.png}\\Verifier $\mathcal{V}^{\textcolor{oc-red-8}{\mathbbm{y}}}(\mathbbm{x})$};

        \draw [dashed,line width=0.3mm] ([yshift=-0.5cm]prover.south) -- ([yshift=-10.4cm]prover.south);
        \draw [dashed,line width=0.3mm] ([yshift=-0.5cm]verifier.south) -- ([yshift=-10.4cm]verifier.south);

        % a strip of cells: #1 = centre, #2 = name, #3 = queried cells, #4 = fill colour
        \newcommand{\strip}[4]{
            \foreach \i in {0,...,9} {
                \node[cell, fill=#4] at ([xshift=-1.1cm+0.36cm*\i]#1) {};
            }
            \foreach \i in {#3} {
                \node[cell, fill=oc-orange-3, draw=oc-orange-8, thick] at ([xshift=-1.1cm+0.36cm*\i]#1) {};
            }
            \node[left] at ([xshift=-1.3cm]#1) {#2 $=$};
        }
        % an oracle sent by the prover: #1 = yshift, #2 = name, #3 = queried cells
        \newcommand{\sendoracle}[3]{
            \draw[send] ([yshift=#1]prover.south) coordinate (l) -- (l-|verifier) coordinate[midway] (m);
            \coordinate (s) at ([yshift=0.35cm]m);
            \strip{s}{#2}{#3}{oc-blue-1}
        }
        % a challenge sent by the verifier: #1 = arrow yshift, #2 = name
        \newcommand{\sendchallenge}[2]{
            \draw[send] ([yshift=#1]verifier.south) coordinate (l) -- (l-|prover) node[midway, above=0mm]{Sample challenge #2};
        }

        % the input oracle y: known to the prover, only queried by the verifier
        \path (prover.south) -- (verifier.south) coordinate[midway] (c);
        \coordinate (y) at ([yshift=-1.25cm]c);
        \node[draw=oc-red-6, dashed, thick, rounded corners=2pt, minimum width=5.6cm, minimum height=1.25cm] at ([yshift=0.2cm]y) {};
        \node[text=oc-red-8, font=\small] at ([yshift=0.55cm]y) {input oracle (already committed)};
        \coordinate (s) at ([yshift=-0.05cm]y);
        \strip{s}{$\textcolor{oc-red-8}{\mathbbm{y}}$}{1,6}{oc-red-1}

        \sendoracle{-2.9cm}{$\pi_0$}{2,7}
        \sendchallenge{-4.1cm}{$\alpha_1$}
        \sendoracle{-5.5cm}{$\pi_1$}{4}
        \node at ([yshift=-6.1cm]c) {$\vdots$};
        \sendchallenge{-7.3cm}{$\alpha_r$}
        \sendoracle{-8.7cm}{$\pi_r$}{0,5,8}

        \node[align=center,fill=white!5,thick,below=9.1cm of verifier]{
        \noindent\rule{4.2cm}{0.8pt}\\
        Query a few cells of $\textcolor{oc-red-8}{\mathbbm{y}}$ and each $\pi_i$\\
        Output $\mathsf{accept}/\mathsf{reject}$ \\
        \noindent\rule{4.2cm}{0.8pt}};
    \end{tikzpicture}%
    }

    \caption{An IOP of proximity. Compared to \Cref{fig:iop-definition}, there is an additional input oracle $\mathbbm{y}$: the prover knows it in full, while the verifier can only query its cells, in the same way as the cells of the oracles $\pi_0, \dots, \pi_r$ (queried cells are marked in orange).}
    \label{fig:iopp-definition}
\end{figure}

Before defining soundness, we need to decide which inputs the verifier is supposed to reject.
The natural answer, ``every $\mathbbm{y}$ that is not valid'', is unfortunately unachievable.
Take a valid $\mathbbm{y}$ and change a single cell: a verifier that queries $q$ of the $n$ cells reads the changed one with probability at most $q/n$, so it almost never notices the difference.
Hence, we only ask the verifier to reject oracles that are wrong in \emph{many} positions, that is, oracles that are far from every valid one.

\begin{definition}
  Let $\mathcal{R}$ be the relation as in \Cref{def:iopp}. 
  We say that $(\mathbbm{x}, \mathbbm{y})$ is \emph{$\delta$-far from $\mathcal{R}$} if 
  \begin{equation*}
    (\mathbbm{x},\mathbbm{y}',\mathbbm{w}) \not\in\mathcal{R} \; \text{for all} \; \mathbbm{y}',\mathbbm{w} \; \text{such that} \; \delta(\mathbbm{y}, \mathbbm{y}') \leq \delta.
  \end{equation*}
\end{definition}

In other words, $(\mathbbm{x}, \mathbbm{y})$ is $\delta$-far from $\mathcal{R}$ if $\mathbbm{y}$ cannot be made valid by changing at most a $\delta$ fraction of its cells.

\begin{definition}
  An IOPP $(\mathcal{P}, \mathcal{V})$ for $\mathcal{R}$ may have the following properties:
  \begin{itemize}
    \item \textbf{Correctness.} For all $(\mathbbm{x}, \mathbbm{y}, \mathbbm{w}) \in \mathcal{R}$ the verifier $\mathcal{V}$ always accepts $\mathcal{P}$.
    \item \textbf{$\delta$-soundness.} For all $(\mathbbm{x}, \mathbbm{y})$ that are $\delta$-far from $\mathcal{R}$,
    \begin{equation*}
      \forall_{\mathcal{P}^*} \; \Pr[\mathcal{V}^{\mathbbm{y},\pi_0,\dots,\pi_r}(\mathbbm{x},\alpha_1,\dots,\alpha_r) = 1] < \negl(\lambda).
    \end{equation*}
  \end{itemize}
\end{definition}

Note that the two properties leave a gap: an oracle $\mathbbm{y}$ that is not valid, but is also not $\delta$-far, may be either accepted or rejected.
This gap is harmless for our purposes.
If the verifier accepts, then $\mathbbm{y}$ is $\delta$-close to some valid $\mathbbm{y}'$, and for a code and a small enough $\delta$ this $\mathbbm{y}'$ is either unique (\Cref{prop:unique-decoding-distance}) or bounded by some constant (that's where our discussion with list decoding from \Cref{subsubsec:list-decoding} becomes useful).
We can therefore treat the prover as committed to $\mathbbm{y}'$, even if a few cells of $\mathbbm{y}$ are corrupted.

\begin{example}\label{ex:iopp}
    Let $\mathbbm{x} = \mathsf{RS}[\mathbb{F}_7, \Omega := \mathbb{F}_7, k := 2]$, and let $\mathcal{R}$ consist of all triples $(\mathbbm{x}, \mathbbm{y}, \mathbbm{w})$ where $\mathbbm{w} = f \in \mathbb{F}_7[T]^{<2}$ and $\mathbbm{y} = \overline{f}$.
    That is, the valid oracles are exactly the codewords of $\mathsf{RS}[\mathbb{F}_7, \Omega, 2]$.
    Consider three oracles:
    \begin{itemize}
        \item $\mathbbm{y}_1 = (3, 5, 0, 2, 4, 6, 1)$ is the codeword of $f = 3 + 2T$, so it is valid, and the verifier always accepts it.
        \item $\mathbbm{y}_2 = (3, 5, 0, 2, 4, 6, \textcolor{oc-red-8}{0})$ differs from $\mathbbm{y}_1$ in one cell. It is not a codeword, but it is $1/7$-close to one, so it is not $\delta$-far for any $\delta \geq 1/7$, and the verifier is allowed to accept it. In this case, we regard the prover as committed to $f = 3 + 2T$.
        \item $\mathbbm{y}_3 = (0, 1, 4, 2, 2, 4, 1)$ consists of the evaluations of $T^2$. A polynomial of degree less than $2$ agrees with $T^2$ in at most $2$ points, so $\mathbbm{y}_3$ differs from every codeword in at least $5$ cells. Hence, it is $\delta$-far from $\mathcal{R}$ for every $\delta < 5/7$ (see \Cref{section:linear-codes-exercises}, Exercise 5); verifier rejects it with high probability.
    \end{itemize}
\end{example}

A special class of IOPPs from \Cref{ex:iopp} where the statement $\mathbbm{x}$ is a code $\code = \mathsf{RS}[\mathbb{F},\Omega,k]$ and $\mathbbm{y}$ is a restriction $g: \Omega \to \mathbb{F}$ is called a \emph{Reed-Solomon IOPP}.
Since this is a central topic of this chapter, we define the RS-IOPP separately:
\begin{definition}[Reed-Solomon IOPP]\label{def:rs-proximity}
  Given oracle access to $g: \Omega \to \mathbb{F}$ and $\delta \in (0,1)$, distinguish, with high confidence and few oracle queries $\pi_0,\dots,\pi_r$, between the case $g \in \mathsf{RS}[\mathbb{F}, \Omega, k]$ and the case that $g$ is $\delta$-far from $\mathsf{RS}[\mathbb{F}, \Omega, k]$.
\end{definition}

\subsubsection{Compiling Poly-IOP into SNARKs using RS-IOPP}

We bet the name of this subsection may cause a lot of confusion, so let us decipher what it means. 
Recall that Poly-IOP is a special class of IOPs where messages of a prover are polynomials of bounded degree while the challenges by a verifier are points at which she wishes to know the value of these polynomials.
To compile Poly-IOP into SNARKs, one employs the following steps:
\begin{enumerate}
  \item First, we replace sent polynomials with \emph{polynomial commitments} to them. 
  For example, the PlonK proving system, considered in \Cref{section:plonk}, is a Poly-IOP with KZG commitments (\Cref{def:kzg-com}); one may for example use Bulletproofs (inner product arguments; see \Cref{section:bulletproofs}) instead to get a Halo2 scheme\footnote{See \url{https://zcash.github.io/halo2/index.html}.}, which we have not discussed yet.
  \item Second, apply the Fiat-Shamir transformation (\Cref{subsection:fiat-shamir}).
\end{enumerate}

Now, the key idea of this section is that \emph{RS-IOPP can be used as a polynomial commitment scheme}, consequently transforming the Poly-IOP into the IOP. 

So now let us think how to do that. First, since IOP requires sending large \emph{strings}, and not \emph{polynomials}, how do we encode a polynomial $f \in \mathbb{F}[T]^{<k}$ from the Poly-IOP into a string of length $n$?
Given the machinery of \Cref{section:linear-codes}, the natural answer is to use the restriction $\overline{f}: \Omega \to \mathbb{F}$ (or alternatively, the evaluation $\overline{f} = (f(\omega))_{\omega \in \Omega}$), which lies in $\code := \mathsf{RS}[\mathbb{F},\Omega,k]$.
The problem, though, that the prover now can send arbitrary functions $g: \Omega \to \mathbb{F}$, so we need to verify that $g$ is ``close enough'' to $\code$.
But that is exactly the instance of RS-IOPP! (\Cref{def:rs-proximity})

However, that is not enough. 
For a polynomial commitment scheme, we need to provide \emph{evaluation proofs}: \emph{i.e.}, that $f(r)=y$ for a given $(r,y)$.
Here is the proposition that can be used.

\begin{proposition}[Quotienting]\label{prop:quotienting}
    Let $r \in \mathbb{F} \setminus \Omega$, $y \in \mathbb{F}$, $f \in \mathbb{F}[T]^{<k}$, $\delta \in [0,1]$.
    Define the \emph{quotient map} $\mathsf{q}$ that maps a restriction $g: \Omega \to \mathbb{F}$ to a restriction $(g(T)-y)/(T-r)$. Then,
    \begin{enumerate}
      \item if $g = \overline{f} \in \mathsf{RS}[\mathbb{F}, \Omega, k]$ and $y=f(r)$, then $\mathsf{q}(\overline{g}) \in \mathsf{RS}[\mathbb{F},\Omega,k-1]$;
      \item suppose for all\footnote{Here, we use a short-hand $\mathsf{List}[g,k,\delta] := \mathsf{List}[g,\mathsf{RS}[\mathbb{F},\Omega,k],\delta]$.} $\overline{h} \in \mathsf{List}[g,k,\delta]$ we have $y \neq h(r)$; then, $\mathsf{q}(\overline{g})$ is $\delta$-far from $\mathsf{RS}[\mathbb{F},\Omega,k-1]$.
    \end{enumerate}
\end{proposition}

\begin{proof}
    Note first that $\mathsf{q}(g)$ is well-defined: since $r \notin \Omega$, the denominator $\omega - r$ is non-zero for every $\omega \in \Omega$.

    \textit{(i)} Since $y = f(r)$, the polynomial $f - y$ has a root at $r$, so it is divisible by $T - r$.
    Hence, $u := (f - y)/(T - r)$ is a polynomial of degree less than $k - 1$.
    For every $\omega \in \Omega$ we have $\mathsf{q}(g)(\omega) = (f(\omega) - y)/(\omega - r) = u(\omega)$, so $\mathsf{q}(g) = \overline{u} \in \mathsf{RS}[\mathbb{F}, \Omega, k-1]$ as was stated.

    \textit{(ii)} We prove the contrapositive.
    Suppose that $\mathsf{q}(g)$ is $\delta$-close to $\mathsf{RS}[\mathbb{F}, \Omega, k-1]$: there is $u \in \mathbb{F}[T]^{<k-1}$ such that $\mathsf{q}(g)(\omega) = u(\omega)$ for all but at most a $\delta$ fraction of the points $\omega \in \Omega$.
    Define $h := u \cdot (T - r) + y$, which is a polynomial of degree less than $k$ with $h(r) = y$.
    At every point $\omega$ where $\mathsf{q}(g)$ and $u$ agree, we have
    \begin{equation*}
        g(\omega) = \mathsf{q}(g)(\omega) \cdot (\omega - r) + y = u(\omega) \cdot (\omega - r) + y = h(\omega).
    \end{equation*}
    Therefore, $\delta(g, \overline{h}) \leq \delta$, that is, $\overline{h} \in \mathsf{List}[g, k, \delta]$, and $h(r) = y$.
    This contradicts the assumption that $y \neq h(r)$ for all $\overline{h} \in \mathsf{List}[g, k, \delta]$.
\end{proof}

\begin{example}
    Let $\mathbb{F} = \Omega = \mathbb{F}_7$, $k = 2$ and $r = 0 \notin \Omega$.
    Take $f = 3 + 2T$, so that $g = \overline{f} = (5, 0, 2, 4, 6, 1)$ and $f(r) = 3$.
    \begin{itemize}
        \item For the correct value $y = 3$, we get $\mathsf{q}(g)(\omega) = (3 + 2\omega - 3)/\omega = 2$ for every $\omega \in \Omega$. Hence, $\mathsf{q}(g) = (2, 2, 2, 2, 2, 2)$ is a constant function, which is a codeword of $\mathsf{RS}[\mathbb{F}_7, \Omega, 1]$.
        \item For the wrong value $y = 4$, we get $\mathsf{q}(g)(\omega) = (2\omega - 1)/\omega = 2 - \omega^{-1}$, that is, $\mathsf{q}(g) = (1, 5, 4, 0, 6, 3)$. All six values are distinct, so $\mathsf{q}(g)$ agrees with any constant function in at most one point and is at relative distance $5/6$ from $\mathsf{RS}[\mathbb{F}_7, \Omega, 1]$.
    \end{itemize}
    This matches the second claim: for a small $\delta$, the only element of $\mathsf{List}[g, 2, \delta]$ is $\overline{f}$ itself, and $f(0) = 3 \neq 4$.
\end{example}

Consequently, we might try to apply \Cref{prop:quotiening} to build the Poly-IOP to IOP reduction.

\textit{Attempt \#1.} This attempt is depicted in \Cref{fig:poly-iop-first-attempt}.
Specifically, upon sending the evaluation $g \gets \overline{f}$, when the verifier wants to know the value of $f$ at specific points $r_1,r_2 \leftarrowS \mathbb{F} \setminus \Omega$, the prover runs the $\delta$-RS-IOPP protocol (\Cref{def:rs-proximity}) over $\mathsf{q}_i(g)$ and verifies that $\mathsf{q}_i(g) \in \mathsf{RS}[\mathbb{F},\Omega,k-1]$.
As a result, what does the verifier $\mathcal{V}$ learn? He learns the following.
\begin{bookquote}
  There are some $\overline{f}_1,\overline{f}_2 \in \mathsf{List}[g,k,\delta]$ s.t. $f_1(r_1)=y_1$ and $f_2(r_2)=y_2$.
\end{bookquote}

In principle, this could be fine, but as we observed in \Cref{subsubsec:list-decoding}, the set $\mathsf{List}[g,k,\delta]$ may contain multiple elements for $\delta \geq \mu(\code)/2$, and so it might happen that $f_1 \neq f_2$.

\begin{figure}[h!]
    \centering
    \scalebox{0.8}{%
    \begin{tikzpicture}[
        cell/.style={draw=oc-gray-7, minimum size=0.36cm, inner sep=0pt},
        vcell/.style={draw=oc-teal-7, densely dashed, fill=oc-teal-1, minimum size=0.36cm, inner sep=0pt},
        send/.style={-{Stealth[length=3mm]}, line width=0.4mm},
        iopp/.style={{Stealth[length=3.5mm]}-{Stealth[length=3.5mm]}, line width=1.2mm, oc-blue-5},
    ]
        \node[inner sep=0pt, align=center] (prover) {\includegraphics[width=1.25cm]{contents/images/common/prover.png}\\Prover $\mathcal{P}(f)$};
        \node[inner sep=0pt, align=center, right=7cm of prover] (verifier) {\includegraphics[width=1.25cm]{contents/images/common/verifier.png}\\Verifier $\mathcal{V}$};

        \draw [dashed,line width=0.3mm] ([yshift=-0.5cm]prover.south) -- ([yshift=-11.2cm]prover.south);
        \draw [dashed,line width=0.3mm] ([yshift=-0.5cm]verifier.south) -- ([yshift=-11.2cm]verifier.south);
        \path (prover.south) -- (verifier.south) coordinate[midway] (c);

        % a strip of cells: #1 = centre, #2 = name, #3 = cell style
        \newcommand{\strip}[3]{
            \foreach \i in {0,...,9} {
                \node[#3] at ([xshift=-1.1cm+0.36cm*\i]#1) {};
            }
            \node[left] at ([xshift=-1.3cm]#1) {#2 $=$};
        }
        % one evaluation query: #1 = index, #2..#5 = yshifts of challenge, answer, virtual oracle, IOPP
        \newcommand{\evalquery}[5]{
            \draw[send] ([yshift=#2]verifier.south) coordinate (l) -- (l-|prover) node[midway, above=0mm]{Sample point $r_{#1} \in \mathbb{F} \setminus \Omega$};
            \draw[send] ([yshift=#3]prover.south) coordinate (l) -- (l-|verifier) node[midway, above=0mm]{Send claimed value $y_{#1}$};
            \coordinate (s) at ([yshift=#4]c);
            \strip{s}{$\mathsf{q}_{#1}(g)$}{vcell}
            %\node[right, font=\small, text=oc-teal-9] at ([xshift=2.35cm]s) {(virtual)};
            \draw[iopp] ([yshift=#5]prover.south) coordinate (l) -- (l-|verifier) node[midway, above=1mm, text=black]{$\delta$-RS-IOPP: is $\mathsf{q}_{#1}(g)$ close to $\mathsf{RS}[\mathbb{F}, \Omega, k-1]$?};
        }

        % the oracle g
        \draw[send] ([yshift=-1.6cm]prover.south) coordinate (l) -- (l-|verifier);
        \coordinate (s) at ([yshift=-1.25cm]c);
        \strip{s}{$g$}{cell, fill=oc-blue-1}
        \node[right] at ([xshift=2.35cm]s) {$: \Omega \to \mathbb{F}$};
        \node[font=\small, text=oc-gray-7] at ([yshift=-1.95cm]c) {honest prover: $g = \overline{f}$};

        \evalquery{1}{-3.3cm}{-4.3cm}{-5.0cm}{-6.3cm}
        \evalquery{2}{-7.6cm}{-8.6cm}{-9.3cm}{-10.6cm}
    \end{tikzpicture}%
    }

    \caption{First attempt at compiling a Poly-IOP into an IOP. Instead of a polynomial $f \in \mathbb{F}[T]^{<k}$, the prover sends an oracle $g: \Omega \to \mathbb{F}$. Each evaluation query $r_i$ is answered with a claimed value $y_i$, followed by an RS-IOPP for the quotient $\mathsf{q}_i(g) := (g - y_i)/(T - r_i)$. Note that the verifier can compute any of quotient's cells from the corresponding cell of $g$. By \Cref{prop:quotienting}, each accepted query is explained by \emph{some} codeword close to $g$, but different queries may be explained by different codewords.}
    \label{fig:poly-iop-first-attempt}
\end{figure}

\textit{Attempt \#2.} So how do we fix this problem? Turns out that fix is rather simple.
Before the procedure illustrate in \Cref{fig:poly-iop-first-attempt}, the verifier will sample a random $r \leftarrowS \mathbb{F} \setminus \Omega$ and an honest prover will respond with some $y \gets f(r)$.
Now we claim the following (note that the statement greatly resembles the classical Schwartz-Zippel lemma~\Cref{lemma:one-sz}):

\begin{proposition}\label{prop:iopp-bad-event}
    Let $\mathsf{Bad}$ be an event ``there exists $\overline{f}_1 \neq \overline{f}_2 \in \mathsf{List}[g,k,\delta]$ such that $f_1(r) = f_2(r) = y$ for a randomly selected $r \leftarrowS \mathbb{F} \setminus \Omega$''.
    Then,
    \begin{equation*}
      \Pr\left[ \mathsf{Bad} \right] \leq \binom{\#\mathsf{List}[g,k,\delta]}{2} \cdot \frac{k-1}{\#\mathbb{F} - \#\Omega}.
    \end{equation*}
\end{proposition}

\begin{proof}
    Fix a pair of distinct codewords $\overline{f_1} \neq \overline{f_2} \in \mathsf{List}[g,k,\delta]$.
    Then $f_1 - f_2$ is a non-zero polynomial of degree less than $k$, so it has at most $k - 1$ roots in $\mathbb{F}$.
    Since $r$ is uniform over the $\#\mathbb{F} - \#\Omega$ elements of $\mathbb{F} \setminus \Omega$, we get (by \Cref{lemma:one-sz})
    \begin{equation*}
        \Pr\left[ f_1(r) = f_2(r) \right] \leq \frac{k-1}{\#\mathbb{F} - \#\Omega}.
    \end{equation*}
    The event $\mathsf{Bad}$ occurs only if $f_1(r) = f_2(r)$ for at least one of the $\binom{\#\mathsf{List}[g,k,\delta]}{2}$ pairs, and the claim follows by the union bound.
\end{proof}

\begin{corollary}\label{cor:iopp-bad-event}
  Adopt notation from \Cref{prop:iopp-bad-event}. 
  If $\delta < 1-\sqrt{1-\mu(\code)}$ (for $\code=\mathsf{RS}[\mathbb{F},\Omega,k]$), then we have
    \begin{equation*}
      \Pr\left[ \mathsf{Bad} \right] \leq \frac{1}{2\varepsilon_{\delta}^2} \cdot \frac{k-1}{\#\mathbb{F} - \#\Omega},
    \end{equation*}
    where the constant $\varepsilon_{\delta}$ does not depend on the size of the field.
\end{corollary}

\begin{proof}
  Simply apply \Cref{th:johnson-bound} to \Cref{prop:iopp-bad-event}. 
  Also use the fact that $\binom{x}{2} \leq \frac{x^2}{2}$ to simplify the expression.
  The constant $\varepsilon_{\delta}$ is given by 
  \begin{equation*}
    \varepsilon_{\delta} = 2\sqrt{1-\mu(\code)}\,(1-\sqrt{1-\mu(\code)}-\delta).
  \end{equation*}
\end{proof}

\begin{remark}
  Note that by \Cref{cor:iopp-bad-event}, the probability of finding two distinct agreeing $f_1$ and $f_2$ is roughly $\mathsf{const}/\#\mathbb{F}$.
  Consequently, for large fields, this probability is negligible.
  In case one applies low-degree fields (like \texttt{M31} where $\#\mathbb{F} = 2^{31}-1$), this procedure can be repeated multiple times to ensure the desired level of soundness.
\end{remark}

Consequently, the pair $(r,y)$ essentially binds the prover to \emph{only one} restriction $\overline{f} \in \mathsf{List}[g,k,\delta]$.
The only problem we have to eliminate is that we do not know to effectively aggregate two checks $f(r)=y$ and $f(r_i)=y_i$.
The following lemma modifies \Cref{prop:quotienting} to support multiple evaluation checks.

\begin{proposition}[Multiple Quotientings]\label{prop:multiple-quotientings}
  Let $(r_i)_{i \in [t]} \subseteq \mathbb{F} \setminus \Omega$, $(y_i)_{i \in [t]} \subseteq \mathbb{F}$, $f: \Omega \to \mathbb{F}$, and $\delta \in [0,1]$.
  Define two polynomials $\nu,\iota \in \mathbb{F}[T]^{\leq t}$ as follows:
  \begin{equation*}
    \nu(T) := \prod_{i \in [t]} (T-r_i), \quad \iota(r_i) = y_i \; \text{for all} \; i \in [t].
  \end{equation*}

  Define the \emph{quotient map} $\mathsf{q}$ that maps a restriction $g: \Omega \to \mathbb{F}$ to a restriction $(g(T)-\iota(T))/\nu(T)$. Then,
  \begin{enumerate}
    \item if $g = \overline{f} \in \mathsf{RS}[\mathbb{F}, \Omega, k]$ and $y_i=f(r_i)$ for all $i \in [t]$, then the quotienting $\mathsf{q}(\overline{g}) \in \mathsf{RS}[\mathbb{F},\Omega,k-t]$;
    \item suppose for all $\overline{h} \in \mathsf{List}[g,k,\delta]$ we have $y_i \neq h(r_i)$ for some $i \in [t]$; then, $\mathsf{q}(\overline{g})$ is $\delta$-far from $\mathsf{RS}[\mathbb{F},\Omega,k-t]$.
  \end{enumerate}
\end{proposition}

\begin{exercise}
  Prove this fact.
\end{exercise}

All in all, the second attempt is depicted in \Cref{fig:poly-iop-second-attempt}.

\begin{figure}[h!]
    \centering
    \scalebox{0.8}{%
    \begin{tikzpicture}[
        cell/.style={draw=oc-gray-7, minimum size=0.36cm, inner sep=0pt},
        vcell/.style={draw=oc-teal-7, densely dashed, fill=oc-teal-1, minimum size=0.36cm, inner sep=0pt},
        send/.style={-{Stealth[length=3mm]}, line width=0.4mm},
        iopp/.style={{Stealth[length=3.5mm]}-{Stealth[length=3.5mm]}, line width=1.2mm, oc-blue-5},
    ]
        \node[inner sep=0pt, align=center] (prover) {\includegraphics[width=1.25cm]{contents/images/common/prover.png}\\Prover $\mathcal{P}(f)$};
        \node[inner sep=0pt, align=center, right=7cm of prover] (verifier) {\includegraphics[width=1.25cm]{contents/images/common/verifier.png}\\Verifier $\mathcal{V}$};
        \coordinate (ps) at (prover.south);
        \coordinate (vs) at (prover.south -| verifier.south);
        \path (ps) -- (vs) coordinate[midway] (c);

        % the new step compared to the first attempt
        \fill[oc-orange-1, rounded corners=3pt] ([xshift=-0.5cm, yshift=-2.5cm]ps) rectangle ([xshift=0.5cm, yshift=-4.65cm]vs);
        \node[font=\small\bfseries, text=oc-orange-9, anchor=north west, rotate=90] at ([xshift=-1.0cm, yshift=-4.55cm]ps) {new: binding};

        \draw [dashed,line width=0.3mm] ([yshift=-0.5cm]ps) -- ([yshift=-10.1cm]ps);
        \draw [dashed,line width=0.3mm] ([yshift=-0.5cm]vs) -- ([yshift=-10.1cm]vs);

        % a strip of cells: #1 = centre, #2 = name, #3 = cell style
        \newcommand{\strip}[3]{
            \foreach \i in {0,...,9} {
                \node[#3] at ([xshift=-1.1cm+0.36cm*\i]#1) {};
            }
            \node[left] at ([xshift=-1.3cm]#1) {#2 $=$};
        }
        \newcommand{\fromverifier}[2]{
            \draw[send] ([yshift=#1]vs) -- ([yshift=#1]ps) node[midway, above=0mm]{#2};
        }
        \newcommand{\fromprover}[2]{
            \draw[send] ([yshift=#1]ps) -- ([yshift=#1]vs) node[midway, above=0mm]{#2};
        }

        % the oracle g
        \draw[send] ([yshift=-1.6cm]ps) -- ([yshift=-1.6cm]vs);
        \coordinate (s) at ([yshift=-1.25cm]c);
        \strip{s}{$g$}{cell, fill=oc-blue-1}
        \node[right] at ([xshift=2.35cm]s) {$: \Omega \to \mathbb{F}$};
        \node[font=\small, text=oc-gray-7] at ([yshift=-1.95cm]c) {honest prover: $g = \overline{f}$};

        % binding step
        \fromverifier{-3.3cm}{Sample random point $r \leftarrowS \mathbb{F} \setminus \Omega$}
        \fromprover{-4.3cm}{Send claimed value $y$}

        % evaluation query
        \fromverifier{-5.7cm}{Sample point $r_1 \in \mathbb{F} \setminus \Omega$}
        \fromprover{-6.7cm}{Send claimed value $y_1$}

        % joint quotient
        \node[font=\small] at ([yshift=-7.3cm]c) {$\nu := (T - r)(T - r_1), \qquad \iota(r) = y, \;\; \iota(r_1) = y_1$};
        \coordinate (s) at ([yshift=-7.95cm]c);
        \strip{s}{$\dfrac{g - \iota}{\nu}$}{vcell}
        %\node[right, font=\small, text=oc-teal-9] at ([xshift=2.35cm]s) {(virtual)};
        \draw[iopp] ([yshift=-9.5cm]ps) -- ([yshift=-9.5cm]vs) node[midway, above=1mm, text=black]{$\delta$-RS-IOPP: is $(g - \iota)/\nu$ close to $\mathsf{RS}[\mathbb{F}, \Omega, k-2]$?};
    \end{tikzpicture}%
    }

    \caption{Second attempt at compiling a Poly-IOP into an IOP. Compared to \Cref{fig:poly-iop-first-attempt}, the verifier first asks for the value of $g$ at a random point $r$ (highlighted). Every later query $r_1$ is then quotiented \emph{together} with $(r, y)$, using the polynomials $\nu$ and $\iota$ from \Cref{prop:multiple-quotientings}. An accepted query is therefore explained by a codeword that also passes through $(r, y)$, and by \Cref{prop:iopp-bad-event} there is only one such codeword, except with small probability.}
    \label{fig:poly-iop-second-attempt}
\end{figure}

%--------------------------------
\subsection{FRI protocol}\label{subsec:fri}
%--------------------------------
Unlike the polynomial commitment schemes from \Cref{section:commitment-schemes}, FRI is not formalized as a static tuple of algorithms (such as $\mathsf{Setup}$ and $\mathsf{Com}$). Instead, it operates as a two-phase interactive protocol, as outlined in \Cref{fig:fri-scheme}.

\begin{figure}[h!]
    \centering
    \scalebox{0.8}{
    \begin{tikzpicture}[
        arr/.style   = {-{Stealth[length=2.5mm]}, line width=0.4mm},
        act/.style   = {draw=oc-blue-7, fill=oc-blue-1, rounded corners=2pt,
                        line width=0.4mm, inner sep=4pt, font=\small},
        data/.style  = {font=\small, inner sep=2pt},
        phase/.style = {draw=gray, dashed, rounded corners=8pt, line width=0.4mm},
        halo/.style  = {preaction={draw, white, line width=2.2mm}},
    ]
        % --- Verifier column
        \draw[fill=gray!20!white, ultra thick, rounded corners=4pt] (8.6,-3.9) rectangle (10.4,6.0);
        \node[inner sep=0pt, align=center] at (9.5,7.0)
            {\includegraphics[width=1.0cm]{contents/images/common/verifier.png}\\ Verifier $\mathcal{V}$};

        % --- Prover
        \node[inner sep=0pt, align=center] at (0,7.0)
            {\includegraphics[width=1.0cm]{contents/images/common/prover.png}\\ Prover $\mathcal{P}$};

        % --- Phase frames
        \draw[phase] (-1.9,0.55) rectangle (8.3,6.0);
        \draw[phase] (-1.9,-3.9) rectangle (8.3,0.15);
        \node[font=\scriptsize\scshape, text=gray] at (6.9,5.6) {Commit phase};
        \node[font=\scriptsize\scshape, text=gray] at (6.9,-0.25) {Query phase};

        % --- Commit phase
        \node[data] (poly)  at (0,5.45) {polynomial $f$};
        \node[act]  (eval)  at (0,4.55) {evaluate on $D_0$};
        \node[data] (cw)    at (0,3.65) {codeword $f_i$ on $D_i$};
        \node[act]  (merk)  at (3.3,2.85) {Merkle commit};
        \node[data] (root)  at (6.3,2.85) {Merkle root};
        \node[act]  (fold)  at (0,2.0) {split-and-fold};
        \node[data] (beta)  at (6.3,2.0) {challenge $\beta_i$};
        \node[data] (final) at (0,1.1) {final value $f_r$};

        \draw[arr] (poly) -- (eval);
        \draw[arr] (eval) -- (cw);
        \draw[arr] (cw.south east) -- (merk.north west);
        \draw[arr] (merk) -- (root);
        \draw[arr] (root) -- (8.6,2.85);
        \draw[arr] (8.6,2.0) -- (beta);
        \draw[arr] (beta) -- (fold);
        \draw[arr] (fold) -- (cw);
        \draw[arr] (fold) -- (final);
        \draw[arr] (final) -- (8.6,1.1);

        % --- Query phase
        \node[data] (trees) at (3.3,-0.7) {Merkle trees};
        \node[act]  (open)  at (3.3,-1.75) {open};
        \node[data] (idx)   at (6.3,-1.75) {query points $s$};
        \node[data] (vals)  at (1.4,-2.8) {values};
        \node[data] (paths) at (5.0,-2.8) {authentication paths};

        \draw[arr, halo] (merk) -- (trees);
        \draw[arr] (trees) -- (open);
        \draw[arr] (8.6,-1.75) -- (idx);
        \draw[arr] (idx) -- (open);
        \draw[arr] (open) -- (vals);
        \draw[arr] (open) -- (paths);
        \draw[arr, rounded corners=4pt] (vals.south) |- (8.6,-3.45);
        \draw[line width=0.4mm, rounded corners=4pt] (paths.south) |- (5.6,-3.45);
    \end{tikzpicture}
    }
    \caption{Overview of the FRI protocol. In the commit phase, the prover repeatedly commits to the
    current codeword and folds it using the verifier's challenge, until a final value remains.
    In the query phase, the verifier asks to open the committed codewords at random points and checks
    the authentication paths. Adapted from~\cite{szepieniec2021anatomy}.}
    \label{fig:fri-scheme}
\end{figure}

To see how the protocol works, we first describe its main building block.

\subsubsection{Split-and-Fold scheme}
FRI follows a split-and-fold strategy: a claim is split into two claims of half the size, the verifier sends a random challenge, and the prover merges the two claims into one using this challenge. Then, after logarithmically many steps, the claim has been reduced to one of trivial size. If the original claim is false, then with high probability the final claim is false as well.

In the FRI case, this claim asserts that the given codeword corresponds to a polynomial of low degree. Let $n$ be the length of the codeword, and let $f(T)=\sum_{i=0}^{k-1} c_i\cdot T^i$ be a polynomial of degree less than $k$.

Following the divide-and-conquer strategy of the fast Fourier transform, we split a polynomial $f$ into its even and odd parts:
\begin{gather*}
  f(T) = f_E(T^2) + T \cdot f_O(T^2),
  \qquad
  f_E(T^2) = \frac{f(T) + f(-T)}{2},
  \\
  f_O(T^2) = \frac{f(T) - f(-T)}{2T}.
\end{gather*}
Indeed, the odd terms cancel in $f(T) + f(-T)$ and the even terms cancel in $f(T) - f(-T)$.
If $\deg f < k$ for an even $k$, then both $f_E$ and $f_O$ have degree less than $k/2$.
Given a random challenge $\beta \in \mathbb{F}_q$ from the verifier, the prover derives a codeword for
\begin{equation*}
  f^\star(T) := f_E(T) + \beta \cdot f_O(T)
\end{equation*}
from the codeword for $f(T)$.

Let $D$ be a subgroup of even order $n$ of the multiplicative group of the field, and let $\omega$ be a generator of this subgroup. Let $\{f(\omega^i)\}_{i=0}^{n-1}$ be the codeword for $f(T)$, corresponding to evaluation on $D$. Let $D^\star=\langle \omega^2\rangle$ be another domain of half the length, and let $\{f_E(\omega^{2i})\}_{i=0}^{n/2-1}$, $\{f_O(\omega^{2i})\}_{i=0}^{n/2-1}$ and $\{f^\star(\omega^{2i})\}_{i=0}^{n/2-1}$ be the codewords for $f_E(T)$, $f_O(T)$ and $f^\star(T)$, respectively, corresponding to evaluation on $D^\star$. As a result, we can rewrite the definition of $f^\star(T)$:
\begin{gather*}
  \{f^\star(\omega^{2i})\}_{i=0}^{n/2-1} = \left\{f_E(\omega^{2i})+\beta\cdot f_O(\omega^{2i})\right\}_{i=0}^{n/2-1} = \\
  \left\{\frac{f(\omega^i)+f(-\omega^i)}{2}+\beta \cdot\frac{f(\omega^i)-f(-\omega^i)}{2\omega^i}\right\}_{i=0}^{n/2-1} = \\
  \left\{2^{-1}\cdot \big((1+\beta\cdot \omega^{-i})\cdot f(\omega^i)+(1-\beta\cdot \omega^{-i})\cdot f(-\omega^i)\big)\right\}_{i=0}^{n/2-1}.
\end{gather*}

Note that, since $\omega^{n/2}=-1$, we have $f(-\omega^i)=f(\omega^{n/2+i})$. This makes it clear that even though the index iterates over half the range, all $n$ points of the codeword for $f(T)$ are involved.

At this point it is possible to describe one round of FRI: the prover commits to $f(T)$ by sending the Merkle root of its codeword to the verifier. The verifier responds with the random challenge $\beta$. The prover computes $f^\star(T)$ and commits to it by sending the Merkle root of $\{f^\star(\omega^{2i})\}_{i=0}^{n/2-1}$ to the verifier.

In turn, having received the two commitments, the verifier can check that they are correctly related, and rejects the proof if for some $x \in D$
\begin{equation}\label{eq:fri-fold}
  f^\star(x^2)\neq 2^{-1}\cdot \big((1+\beta x^{-1})\cdot f(x)+(1-\beta x^{-1})\cdot f(-x)\big).
\end{equation}
To do this, the verifier randomly samples an index $i \leftarrowS [n/2]$, sets $x := \omega^i$, and thus defines $3$ points:
\begin{itemize}
  \item $A:\quad (\omega^i, f(\omega^i))$,
  \item $B:\quad (\omega^{n/2+i}, f(\omega^{n/2+i}))$,
  \item $C:\quad (\beta, f^\star(\omega^{2i}))$.
\end{itemize}
Upon receiving the index $i$ from the verifier, the prover provides the $y$-coordinates along with their Merkle authentication paths. The verifier checks these paths against the corresponding roots and then verifies that $A$, $B$ and $C$ lie on a straight line. This is called the \emph{collinearity check}.

%Add exercise A and B on the same line.

After the first round is complete, the verifier and the prover return to the starting point. The prover wishes to establish that a given Merkle root decommits to a codeword whose defining polynomial has a bounded degree. However, as a result of the previous round, the length of the codeword and the number of possible non-zero coefficients of the polynomial have both been halved. Now they set $f := f^\star$ and $D := D^\star$ and repeat the process with a fresh challenge. If $k = 2^r$, then after $r = \log_2 k$ rounds of FRI, the prover and the verifier end up with a constant polynomial whose codeword is also constant, so the prover simply sends this constant to the verifier.

\begin{definition}[FRI protocol]\label{def:fri}
  Let $f_0 := f$ be given on $D_0 := D = \langle \omega \rangle$ of order $n$, with $\deg f < k = 2^r$.
  \begin{itemize}
    \item \textbf{Commit phase.} For $i = 0, 1, \ldots, r-1$:
    \begin{algoen}
      \item the prover sends the Merkle root of the codeword of $f_i$ on $D_i$;
      \item the verifier sends a random challenge $\beta_i \leftarrowS \mathbb{F}_q$;
      \item the prover computes $f_{i+1} := f_{i,E} + \beta_i \cdot f_{i,O}$ on $D_{i+1} := \{x^2 : x \in D_i\}$.
    \end{algoen}
    Finally, the prover sends the constant $f_r$.
    \item \textbf{Query phase.} Repeat $t$ times:
    \begin{algoen}
      \item the verifier samples $x_0 \leftarrowS D_0$ and sets $x_{i+1} := x_i^2$ for $i \in [r]$;
      \item for every $i \in [r]$, the prover opens $f_i(x_i)$ and $f_i(-x_i)$ together with their
      authentication paths;
      \item the verifier checks all paths and \eqref{eq:fri-fold} for every $i \in [r]$, where
      $f_r(x_r)$ is the constant sent by the prover.
    \end{algoen}
  \end{itemize}
  The verifier accepts if and only if all checks pass.
\end{definition}

\begin{remark}
  As described, FRI is interactive. In STARKs it is made non-interactive using the Fiat-Shamir
  transformation (\Cref{subsection:fiat-shamir}): each challenge $\beta_i$ and each starting point $x_0$
  is derived by hashing the Merkle roots sent so far.
\end{remark}

\subsection{Arithmetic Intermediate Representation (AIR)}\label{subsec:air}
Similar to SNARK, the STARK proof system involves intermediate stages, the first two of which are the \emph{arithmetic constraint system} and the \emph{IOP}.

\begin{definition}\label{def:air}
  The \emph{arithmetic intermediate representation (AIR)} or arithmetic internal representation, is a way of describing a computation in terms of an execution trace that satisfies a number of constraints induced by the correct evolution of the state.
\end{definition}

\begin{definition}\label{def:aet}
  Let $\mathbb{F}_q$ be the field of definition, and the computation describes the evolution of a state of $w$ registers for $\tau$ cycles. Then we define the \emph{algebraic execution trace (AET)} is the table of $\tau \times w$ field elements where every row describes the state of the system at the given point in time, and every column tracks the value of the given register.
\end{definition}

\begin{definition}\label{def:st-function}
  A \emph{state transition function} $\phi:\mathbb{F}_q^{w}\rightarrow\mathbb{F}_q^{w}$ determines the state at the next cycle as a function of the state at the previous cycle. Furthermore, a set of boundary conditions 
  \begin{equation*}
    \mathcal{B}\subseteq \mathbb{Z}_\tau \times\mathbb{Z}_w\times\mathbb{F}_q
  \end{equation*}
  fixes the values of some registers at some cycles: a triple $(i,j,e) \in \mathcal{B}$ means that register $j$ holds the value $e$ at cycle $i$.
\end{definition}

The \emph{computational integrity claim} consists of the state transition function and the boundary conditions. The \emph{witness} to this claim is the algebraic execution trace. The claim is \emph{true} if there is a witness $W\in \mathbb{F}_q^{\tau \times w}$ such that:
\begin{itemize}
  \item For every cycle, the state evolves correctly: $\forall i\in [\tau-1], \phi(W_{[i, :]})=W_{[i+1, :]}$
  \item All boundary conditions are satisfied: $\forall (i, j, e)\in \mathcal{B}, W_{[i, j]}=e$
\end{itemize}
Actually, the state transition function hides a lot of complexity and requires to be describable as low degree polynomials that are independent of the cycle.
However, this list of polynomials does not need to compute the next state from the current one; it merely needs to distinguish correct evolutions from incorrect ones. Specifically, the state transition function $\phi$ is represented by a list of polynomials $\mathbf{p}(x_0,\ldots,x_{w-1}, y_0,\ldots,y_{w-1})$ such that $\phi(\mathbf{x})=\mathbf{y}$ if and only if $\mathbf{p}(\mathbf{x},\mathbf{y})=0$. Say there are $m$ such state transition verification polynomials. Then the transition constraint become:
\begin{equation*}
  \forall i\in[\tau-1],\forall  \ell\in[m]: p_\ell(W_{[i, 0]},\ldots,W_{[i, w-1]},W_{[i+1, 0]},\ldots,W_{[i+1, w-1]})=0
\end{equation*}
This representation admits non-determinism: the verification polynomials may have much lower degree than any polynomial that computes the next state. For instance, the transition $y = x^{-1}$ requires computing $x^{q-2}$, but it is verified by the degree-$2$ polynomial $xy - 1$.

\subsubsection{Trace interpolation}
We saw that the arithmetic constraint system described above represents the computational integrity claim as a bunch of polynomials, where each such polynomial corresponds to a constraint. However, we need to extend terms of polynomials to the witness for transformation this constraint system into an IOP, and extending the notion of valid witnesses to witness polynomials.

Let $H = \langle h \rangle$ be a multiplicative subgroup of $\mathbb{F}_q^{\times}$ of order $\tau$, which we call the \emph{trace evaluation domain}. For simplicity, we assume that $\tau$ is a power of two; otherwise the trace is padded. The trace domain is smaller than the FRI evaluation domain $D$ by a factor exactly equal to the blowup factor $1/\rho$.

For every register $j \in [w]$, the \emph{trace polynomial} $t_j \in \mathbb{F}_q[T]$ is the unique polynomial of degree less than $\tau$ such that
\begin{equation*}
  t_j(h^i) = W_{[i,j]} \quad \forall i \in [\tau].
\end{equation*}
In practice, it is computed with the inverse NTT (\Cref{section:ntt}). The trace polynomials
$t_0, \ldots, t_{w-1}$ represent the algebraic execution trace in terms of univariate polynomials.

Translating the computational integrity claim to the trace polynomials, we get:
\begin{itemize}
  \item All boundary constraints are satisfied: $\forall (i,j,e)\in \mathcal{B}, t_j(h^i)=e$.
  \item For all cycles, all transition constraints are satisfied: $\forall i\in [\tau-1],\forall \ell\in[m]$
  \begin{equation*}
    p_\ell(t_0(h^i),\dots,t_{w-1}(h^i),t_0(h^{i+1}),\dots,t_{w-1}(h^{i+1}))=0
  \end{equation*}
\end{itemize}
If we look at the left-hand side of the equation, we can see that it corresponds to a univariate polynomial $p_\ell(t_0(T),\dots,t_{w-1}(T),t_0(h\cdot T),\dots,t_{w-1}(h\cdot T))$. Which just means that all $m$ of these transition polynomials evaluate to $0$ in $H$.

Then we can build the following high-level protocol:
This observation gives rise to the following high-level protocol:
\begin{algoen}
  \item The prover commits to the trace polynomials $t_0(T), \ldots, t_{w-1}(T)$.
  \item The verifier checks that $t_j(h^i) = e$ for all $(i,j,e) \in \mathcal{B}$.
  \item The prover commits to the transition polynomials
  \begin{equation*}
    c_\ell(T) := p_\ell\big(t_0(T), \ldots, t_{w-1}(T), t_0(h \cdot T), \ldots, t_{w-1}(h \cdot T)\big),
    \quad \ell \in [m].
  \end{equation*}
  \item The verifier checks that $c_\ell(T)$ and $t_j(T)$ are correctly related by:
  \begin{itemize}
    \item choosing a random point $z \leftarrowS \mathbb{F}_q^{\times}$;
    \item querying the values $t_j(z)$ and $t_j(h \cdot z)$ $\forall j \in [w]$;
    \item evaluating the transition verification polynomials $p_\ell$ at
    \begin{equation*}
      \big(t_0(z), \ldots, t_{w-1}(z), t_0(h \cdot z), \ldots, t_{w-1}(h \cdot z)\big)
    \end{equation*}
    \item querying the values $c_\ell(z)$;
    \item checking that the values obtained in the previous two steps match.
  \end{itemize}
  \item The verifier checks that the transition polynomials $c_\ell(T)$ evaluate to zero on
  $\{h^i : i \in [\tau-1]\}$.
\end{algoen}

In fact, the commitment to the transition polynomials can be omitted. Instead, the verifier uses
the values of $t_j$ at $z$ and $h \cdot z$ to compute the value of $c_\ell(T)$ at the one point
needed to verify that $c_\ell(T)$ evaluates to $0$ on $\{h^i : i \in [\tau-1]\}$.

We could use the protocol above if it did not have a flaw: the verifier cannot perform steps 2 and 5, because he sees the values of the polynomials only on the domain $D$, while these checks concern the trace domain $H$. That is why we use FRI to simulate these evaluation checks, with a trick similar to the one used for KZG polynomial commitments. By \Cref{th:root-div}, $g(x) = y$ if and only if $(T - x) \mid g(T) - y$. Hence an evaluation check can be simulated by:
\begin{itemize}
  \item subtracting the $y$-coordinate;
  \item dividing out the \emph{zerofier}, which is the minimal polynomial that vanishes at the $x$-coordinate;
  \item proving with FRI that the resulting quotient has a bounded degree.
\end{itemize}

This procedure happens twice for the STARK polynomials: (1) first, applied to the trace polynomials to show satisfaction of the boundary constraints and (2) second, applied to the transition polynomials to show that the transition constraints are satisfied. We call the resulting lists of quotient polynomials the \emph{boundary quotients} and the \emph{transition quotients} respectively.

Both kinds of constraints are instances of a single, more general notion.

\begin{definition}[Generalized AIR constraint]\label{def:gen-air}
  A \emph{generalized AIR constraint} is given by two polynomials:
  \begin{itemize}
    \item The \emph{numerator} determines which equation between elements of the algebraic execution
    trace holds, in a manner that is independent of the cycle. For transition constraints, the numerator
    is exactly the transition polynomial $c_\ell(T)$. For a boundary constraint $(i,j,e) \in \mathcal{B}$,
    the numerator is simply $t_j(T) - e$.
    \item The \emph{denominator} is a zerofier that determines where the equation is supposed to hold,
    by vanishing in those points. For transition constraints, this zerofier vanishes on all points of the
    trace evaluation domain except the last. For boundary constraints, it vanishes only on the boundary points.
  \end{itemize}
  The constraint is satisfied if and only if the numerator is divisible by the denominator, that is,
  if the \emph{quotient} of the numerator by the denominator is a polynomial.
\end{definition}

For every register $j \in [w]$, let $\mathcal{B}_j := \{h^i : (i,j,e) \in \mathcal{B}\}$ be its set of
boundary points, let the \emph{boundary interpolant} $I_j$ be the polynomial of degree less than
$|\mathcal{B}_j|$ such that $I_j(h^i) = e$ for all $(i,j,e) \in \mathcal{B}$, and let
$Z_{\mathcal{B}_j}(T) := \prod_{x \in \mathcal{B}_j} (T - x)$ be the \emph{boundary zerofier}.
The \emph{transition zerofier} is
\begin{equation*}
  Z(T) := \prod_{i \in [\tau-1]} (T - h^i) = \frac{T^\tau - 1}{T - h^{\tau-1}},
\end{equation*}
where the second equality follows from \Cref{lemma:vanishing-polynomial}. The quotients are
\begin{equation*}
  q^{\mathcal{B}}_j(T) := \frac{t_j(T) - I_j(T)}{Z_{\mathcal{B}_j}(T)}, \quad j \in [w],
  \qquad
  q^{\mathcal{T}}_\ell(T) := \frac{c_\ell(T)}{Z(T)}, \quad \ell \in [m].
\end{equation*}

The degree bounds are $\deg q^{\mathcal{B}}_j < \tau - |\mathcal{B}_j|$ and, if $p_\ell$ has degree $d_\ell$, $\deg q^{\mathcal{T}}_\ell \leq (d_\ell - 1)(\tau - 1)$. As a result, we can summarize the discussion as follows.

\begin{proposition}\label{prop:stark-quotients}
  The trace $W$ satisfies all boundary and transition constraints if and only if all boundary quotients $q^{\mathcal{B}}_j$ and all transition quotients $q^{\mathcal{T}}_\ell$ are polynomials.
\end{proposition}

\begin{proof}
  A quotient is a polynomial if and only if its numerator vanishes at every root of its denominator (\Cref{th:root-div}, applied to each root in turn). For the boundary quotients this means $t_j(h^i) = e$ for all $(i,j,e) \in \mathcal{B}$, and for the transition quotients it means $c_\ell(h^i) = 0$ for all $i \in [\tau-1]$, which are exactly the constraints.
\end{proof}

The trace polynomial $t_j$ appears in both types of quotients, while $I_j$ and $Z_{\mathcal{B}_j}$ are public. Therefore, the trace can be reconstructed from the boundary quotient,
\begin{equation*}
  t_j(T) = q^{\mathcal{B}}_j(T) \cdot Z_{\mathcal{B}_j}(T) + I_j(T),
\end{equation*}
and the prover commits to the boundary quotients instead of the trace polynomials. After opening $q^{\mathcal{B}}_j$ at the points $x$ and $h \cdot x$, the verifier computes $t_j$ itself, then derives $c_\ell$ from it, and finally derives $q^{\mathcal{T}}_\ell$ from $c_\ell$.

\begin{remark}
  The zerofiers $Z$ and $Z_{\mathcal{B}_j}$ vanish on points of $H$, so the quotients cannot be evaluated there. Since $H = \langle h \rangle$ is a subgroup of $\langle \omega \rangle$, we take the FRI domain to be a coset $D = c \cdot \langle \omega \rangle$ with $c \notin \langle \omega \rangle$, which does not intersect $H$. Moreover, $h = \omega^{1/\rho}$, so $h \cdot x \in D$ for every $x \in D$. FRI works on a coset in exactly the same way: $-x \in D$ for every $x \in D$, and since $(cx)^2 = c^2 x^2$, squaring maps $D$ onto the coset $c^2 \cdot \langle \omega^2 \rangle$ of half the size.
\end{remark}

\subsubsection{Nonlinear Combination}

Instead of running FRI separately for every quotient, the prover combines all of them into a single polynomial. For simplicity, we assume from now on that all transition verification polynomials have degree at most $2$, as in~\cite{szepieniec2021anatomy}. Then every quotient has degree less than $\tau$, and we run FRI with the degree bound $k := \tau$, so that $n = \tau/\rho$.

Let $\mathcal{Q}$ denote the list of all boundary and transition quotients, and for $q \in \mathcal{Q}$ let $d_q$ be its degree bound: $d_q = \tau - |\mathcal{B}_j| - 1$ for $q = q^{\mathcal{B}}_j$, and $d_q = (d_\ell - 1)(\tau - 1)$ for $q = q^{\mathcal{T}}_\ell$. After receiving random weights $\alpha_q, \beta_q \leftarrowS \mathbb{F}_q$ from the verifier, the prover computes
\begin{equation*}
  g(T) := \sum_{q \in \mathcal{Q}} \big( \alpha_q \cdot q(T) + \beta_q \cdot T^{\,k-1-d_q} \cdot q(T) \big).
\end{equation*}

Each quotient enters the sum twice. The first term alone would not be enough: a quotient whose degree exceeds its own bound $d_q$ but is still less than $k$ would not affect the degree of $g$ and would go unnoticed. The second term shifts $q$ so that its degree bound becomes exactly $k-1$: if $\deg q > d_q$, this term has degree at least $k$. Hence, with high probability over the choice of the weights, $\deg g < k$ if and only if $\deg q \leq d_q$ for every $q \in \mathcal{Q}$. Likewise, if some quotient is not a polynomial at all, then with high probability the codeword of $g$ on $D$ is far from $\mathsf{RS}_q[D,k]$. Following~\cite{szepieniec2021anatomy}, we call $g$ a \emph{nonlinear combination}, because of the factors $T^{\,k-1-d_q}$. It therefore suffices to run FRI once, on $g$.

\subsection{The STARK Protocol}\label{subsec:stark-protocol}

We can now put everything together. Note that the random point $z$ from the high-level protocol is no
longer needed: its role is played by the points that the verifier queries in the query phase of FRI.

\begin{definition}[STARK protocol]\label{def:stark}
  The AIR $\big(\tau, w, \{p_\ell\}_{\ell \in [m]}, \mathcal{B}\big)$, the domains $H$ and $D$, and the
  degree bound $k$ are public; the prover knows a trace $W$.
  \begin{algoen}
    \item The prover computes the trace polynomials $t_j$ and the boundary quotients $q^{\mathcal{B}}_j$,
    and sends the Merkle roots of the codewords of $q^{\mathcal{B}}_j$ on $D$ for all $j \in [w]$.
    \item The verifier sends random weights $\alpha_q, \beta_q \leftarrowS \mathbb{F}_q$ for all
    $q \in \mathcal{Q}$.
    \item The prover computes the transition quotients $q^{\mathcal{T}}_\ell$ and the nonlinear
    combination $g$, and runs the commit phase of FRI (\Cref{def:fri}) with $f_0 := g$.
    \item The parties run the query phase of FRI. For every point $x \in D$ at which the verifier reads
    $f_0$, the prover additionally opens $q^{\mathcal{B}}_j(x)$ and $q^{\mathcal{B}}_j(h \cdot x)$ for all
    $j \in [w]$, together with their authentication paths.
    \item For every such $x$, the verifier checks the authentication paths, recovers $t_j(x)$ and
    $t_j(h \cdot x)$ from the boundary quotients, computes $c_\ell(x)$,
    $q^{\mathcal{T}}_\ell(x) = c_\ell(x)/Z(x)$ and $g(x)$, and checks that $g(x) = f_0(x)$.
  \end{algoen}
  The verifier accepts if and only if all these checks and all checks of FRI pass.
\end{definition}

\begin{figure}[h!]
    \centering
    \scalebox{0.74}{
    \begin{tikzpicture}[
        arr/.style   = {-{Stealth[length=2.5mm]}, line width=0.4mm},
        act/.style   = {draw=oc-blue-7, fill=oc-blue-1, rounded corners=2pt,
                        line width=0.4mm, inner sep=4pt, font=\small},
        data/.style  = {draw=gray, rounded corners=2pt, line width=0.4mm,
                        inner sep=4pt, font=\small},
        lbl/.style   = {font=\scriptsize, inner sep=2pt},
        phase/.style = {draw=gray, dashed, rounded corners=8pt, line width=0.4mm},
        proc/.style  = {draw=gray, fill=gray!20!white, rounded corners=2pt,
                        line width=0.4mm, inner sep=4pt, font=\small},
        check/.style = {draw=black, rounded corners=2pt, line width=0.6mm,
                        inner sep=4pt, font=\small},
    ]
        % ---------------- Prover row
        \draw[phase] (-0.8,1.05) rectangle (13.1,-2.3);
        \node[font=\scriptsize\scshape, text=gray, anchor=north east] at (13.0,0.95) {Prover};
        \node[inner sep=0pt, align=center] at (-1.6,-0.6)
            {\includegraphics[width=1.0cm]{contents/images/common/prover.png}\\ $\mathcal{P}$};

        \node[data] (W)  at (-0.15,0)   {$W$};
        \node[data] (tj) at (2.3,0)     {$t_j$};
        \node[act]  (qB) at (6.7,0)     {$q^{\mathcal{B}}_j$};
        \node[data] (cl) at (4.6,-1.55) {$c_\ell$};
        \node[data] (qT) at (6.7,-1.55) {$q^{\mathcal{T}}_\ell$};
        \node[act]  (g)  at (9.2,-0.78) {$g$};
        \node[proc] (fri) at (11.8,-0.78) {FRI (\Cref{fig:fri-scheme})};

        \draw[arr] (W) -- node[lbl, above] {inverse NTT} (tj);
        \draw[arr] (tj) -- node[lbl, above] {$(t_j - I_j)/Z_{\mathcal{B}_j}$} (qB);
        \draw[arr, rounded corners=4pt] (tj.south) |- node[lbl, pos=0.75, below] {$p_\ell$} (cl.west);
        \draw[arr] (cl) -- node[lbl, above] {$\div Z$} (qT);
        \draw[arr] (qB.east) -- (g.north west);
        \draw[arr] (qT.east) -- (g.south west);
        \draw[arr] (g) -- (fri);
        \node[lbl, above=1mm of g] {$\alpha_q, \beta_q$};

        % ---------------- Verifier row
        \draw[phase] (-0.8,-2.75) rectangle (13.1,-5.35);
        \node[font=\scriptsize\scshape, text=gray, anchor=north east] at (13.0,-2.85)
            {Verifier, at every point $x$ queried by FRI};
        \node[inner sep=0pt, align=center] at (-1.6,-4.05)
            {\includegraphics[width=1.0cm]{contents/images/common/verifier.png}\\ $\mathcal{V}$};

        \node[data] (v1) at (1.05,-3.9) {$q^{\mathcal{B}}_j(x),\ q^{\mathcal{B}}_j(h x)$};
        \node[data] (v2) at (4.6,-3.9)  {$t_j(x),\ t_j(h x)$};
        \node[data] (v3) at (6.9,-3.9)  {$c_\ell(x)$};
        \node[data] (v4) at (8.8,-3.9)  {$q^{\mathcal{T}}_\ell(x)$};
        \node[check] (v5) at (11.5,-3.9) {$g(x) \stackrel{?}{=} f_0(x)$};

        \draw[arr] (v1) -- (v2);
        \draw[arr] (v2) -- (v3);
        \draw[arr] (v3) -- (v4);
        \draw[arr] (v4) -- (v5);
        \draw[arr, rounded corners=4pt] (v1.south) -- ++(0,-0.55) -| (v5.south);
    \end{tikzpicture}
    }
    \caption{Overview of the STARK protocol. The prover commits only to the highlighted polynomials:
    the boundary quotients $q^{\mathcal{B}}_j$ and, through FRI, the nonlinear combination $g$. The trace
    polynomials $t_j$, the transition polynomials $c_\ell$ and the transition quotients
    $q^{\mathcal{T}}_\ell$ are never committed. At every point queried by FRI, the verifier recomputes
    them from the opened values of the boundary quotients and compares the result with the first layer
    of FRI. Adapted from~\cite{szepieniec2021anatomy}.}
    \label{fig:stark-scheme}
\end{figure}

If the trace is correct, the honest prover is always accepted. If it is not, then by
\Cref{prop:stark-quotients} at least one quotient is not a polynomial of bounded degree, so with high
probability the codeword of $g$ is far from $\mathsf{RS}_q[D,k]$, and FRI rejects. The prover cannot
avoid this by running FRI on a different function that is close to a codeword: step 5 ties the values
of $f_0$ at the queried points to the committed boundary quotients.

\begin{remark}
  As with FRI, the whole protocol is made non-interactive using the Fiat--Shamir transformation
  (\Cref{subsection:fiat-shamir}). Practical STARKs also add several details that we omit: randomization
  of the trace for zero-knowledge, and additional checks that improve soundness, such as DEEP.
\end{remark}

\end{document}